\section{Hodge theory and formal lifting}\label{sec:hodge}

Our goal is to lift functions and a transverse parameter through
every infinitesimal neighborhood of the divisor constructed in
Proposition~\ref{prop:signed-construction}.  We recall its data.
\(\mathcal P=\sqrt[\ell]{\OO_{\PP^s}(1)}\) is a root gerbe, \(C\) is
its tautological line bundle, and
\[
 S\xrightarrow{g}\mathcal Y\xrightarrow{p}\mathcal P,
 \qquad h=p\circ g.
\]
Here \(g\) is representable and projective, and \(p\) is a smooth
projective bundle.  The Stein factor of \(h\) is finite over
\(\mathcal P\); it has been embedded as \(\mathcal T\subset\mathcal Y\)
so that \(g_*\OO_S=\OO_{\mathcal T}\).  The reduced Cartier
divisor \(S\subset\mathcal X\) has dimension \(N=n-1\).  If
\(J=(S\cap T)_{\mathrm{red}}\), then \(S,J\) are Du Bois,
\(\mathcal I_{J/S}\) is maximal Cohen--Macaulay, and
\begin{equation}\label{hg:geometric-lines}
 \mathcal B:=
 \mathcal Hom_S(\mathcal I_{J/S},\omega_S)
 \simeq \OO_S(aS)=g^*C^a,
 \qquad
 \OO_S(S)=g^*C,
\end{equation}
for an integer \(a>0\).  Off \(T\), the fixed ambient isomorphism
\(\omega_{\mathcal X}(S)\simeq\OO_{\mathcal X}(aS)\) supplies these
identifications by adjunction.

Put \(I=\OO_{\mathcal X}(-S)\). For a separated quasi-projective
scheme etale chart \(U\to\mathcal Y\) of finite type, write
\(S_U=S\times_{\mathcal Y}U\) and again \(g:S_U\to U\)
for the induced projective morphism. Denote by \(qS_U\) the
unique extension of \(S_U\to S\) to the nilpotent thickening
\(qS\). These extensions are compatible in \(q\).
Proposition~\ref{prop:signed-construction} makes them quasi-projective
schemes: they are separated and quasi-finite over \(X\), and hence
open in schemes finite over the projective variety \(X\).
Powers and quotients of \(I\) below are pulled back to the appropriate
thickenings. On these quotients, \(g_*\) denotes direct image for
the common underlying topological map; a ringed-space map from a
thickening to \(U\) has not yet been constructed.

\begin{proposition}\label{prop:formal-lifting}
For every such chart \(U\), every \(j\geq0\), and every \(k\geq1\),
the transition
\[
 g_*(I^j/I^{j+k+1})
       \longrightarrow g_*(I^j/I^{j+k})
\]
is surjective as a map of sheaves of abelian groups.  Its kernel is
\(\OO_{\mathcal T_U}\otimes C^{-j-k}\), where
\(\mathcal T_U=\mathcal T\times_{\mathcal Y}U\).
The statements are compatible with further etale changes of
charts.  On an affine chart, these transitions are also
surjective on global sections.
\end{proposition}

In their isolated-component case with numerical dimension one,
Liu--Xu use finite covers and Du Bois cohomology to obtain infinitesimal
motion \cite[Theorem~3.1]{LiuXu2025}. Here the semiample boundary image
may have positive dimension, so its local functions and the transverse
parameter must be lifted together.

For \(j\ge0\) and \(k\ge1\), the kernel of the transition from length \(k+1\) to length \(k\)
in \(I^j/I^{j+k+1}\) is \(\OO_S\otimes C^{-j-k}\).
The corresponding connecting homomorphism therefore takes values in
\[
 R^1g_*\OO_S\otimes C^{-j-k}.
\]
We will identify this sheaf with a negative twist of the lowest
filtered piece \(F_0M\) of a Hodge-module direct image.  Lifting local
functions gives Cech cocycles in this coherent sheaf.  A calculation
with their graphs shows that the coordinate cocycle lies in the kernel
of the differential-operator symbol map.  Negative-twist vanishing
will make that kernel have no global sections.  The same graph
calculation then kills the error in lifting the transverse parameter.

\subsection{The lowest Hodge piece}

We work with graded-polarizable mixed Hodge modules on ordinary
complex algebraic varieties.  The inputs are projective strictness
and the usual functorial operations
\cite[Theorem~2.14 and Section~4]{Saito1990}, the comparison with the
Du Bois complex and its coherent dual
\cite[Theorem~0.2, Corollary~0.3 and Proposition~5.3]{Saito2000}, and
Kodaira--Saito vanishing \cite[Proposition~2.33]{Saito1990}.
On a smooth stack, we use compatible systems of these objects on
scheme etale charts and descend their filtered differential modules.
Thus the argument requires no Hodge-module theory on stacks.

All differential modules in this section are right modules with an
increasing filtration.  In the Spencer de Rham complex the module
itself occupies degree zero.  Accordingly, on a smooth chart \(V\),
the term in
degree \(-i\) of \(\Gr^F_k\DR_V(M)\) is
\[
 \Gr^F_{k-i}M\otimes_{\OO_V}\bigwedge^iT_V.
\]
When \(F_{<0}M=0\), the lowest graded de Rham complex therefore
consists of the sheaf \(F_0M\) in degree zero.

We now identify a Hodge module whose lowest piece contains the
coherent lifting obstructions. On each chart \(U\), write
\(J_U=J\times_{\mathcal Y}U\).
For \(j:S_U\setminus J_U\hookrightarrow S_U\), put
\[
 A^\bullet_U=\mathbf D\!\left(
             j_!\Q^H_{S_U\setminus J_U}[N]\right),
 \qquad
 A_{0,U}={}^{p}\!H^0A^\bullet_U,
 \qquad
 M_U={}^{p}\!H^1g_*A_{0,U}.
\]
Here \(\mathbf D\) is Hodge-module duality, \({}^pH^i\) denotes
perverse cohomology, and \(g_*\) in the definition of \(M_U\)
is the derived Hodge-module direct image.  Later, \(g_+\) will
denote the corresponding direct image of differential modules.
These constructions commute with etale restriction.  Thus the
\(M_U\) form a system \(M\) supported on \(\mathcal T\).

\begin{lemma}\label{hg:lowest-piece}
The following identifications hold compatibly on the charts:
\[
 F_{<0}A_{0,U}=0,\qquad F_0A_{0,U}=\mathcal B|_{S_U},
 \qquad
 F_{<0}M_U=0,\qquad
 F_0M_U=R^1g_*(\mathcal B|_{S_U}).
\]
In a smooth ambient embedding \(S_U\hookrightarrow V\) of
codimension \(c\), the underlying module of \(A_{0,U}\) is
\(\mathcal H^c_{S_U}(\omega_V)\), localized off \(J_U\).
The lowest-piece inclusion is the adjunction, or Ext-to-support,
inclusion off \(J_U\), extended by meromorphic localization.
\end{lemma}

\begin{proof}
The proof must identify not just a coherent sheaf but its actual
inclusion into the differential module: the graph calculation
will differentiate classes in that inclusion.  We first compute
the sheaf by duality.  Du Bois
comparison applied to the difference triangle for the constant
objects of \(S_U\) and \(J_U\) identifies its degree-zero
graded de Rham complex with \(\mathcal I_{J_U/S_U}\).
Coherent duality and the maximal-Cohen--Macaulay property give
\begin{equation}\label{hg:derived-lowest}
 \Gr^F_k\DR(A^\bullet_U)=0\quad(k<0),
 \qquad
 \Gr^F_0\DR(A^\bullet_U)=\mathcal B|_{S_U}[0].
\end{equation}
Here \([N]\) cancels the dimension shift of the dualizing
complex.  In the smooth case, our convention gives
\(\mathbf D(\Q^H[N])=\Q^H[N](N)\); its right differential module
has \(\omega\) as its lowest piece, at index zero.

To pass to perverse cohomology, work locally on a fixed chart
and let \(q\) be the smallest filtration index occurring in any
perverse cohomology object of \(A^\bullet_U\).  Boundedness of
the complex and goodness of the filtrations provide such a lower
bound.  At index \(q\), each of these objects has a graded
Spencer complex consisting only of its degree-zero term.  The
spectral sequence for perverse truncation has a single nonzero
row at that index, and therefore gives
\[
 \mathcal H^i\Gr^F_q\DR(A^\bullet_U)
   =F_q\,{}^p\!H^iA^\bullet_U.
\]
If \(q<0\), this contradicts Equation~\eqref{hg:derived-lowest}.
It follows that all these objects have \(F_{<0}=0\).  The same
argument at index zero yields
\(F_0A_{0,U}=\mathcal B|_{S_U}\), while the lowest pieces of
the other perverse cohomology objects vanish.  This passage
does not require the shifted constant complex to be perverse.

We next identify the underlying unfiltered module.  Ambient duality in \(V\)
identifies the dual of the reduced constant complex with
\[
 R\Gamma_{[S_U]}(\omega_V)[c].
\]
The degree-zero module is the first nonzero support cohomology,
\(\mathcal H^c_{S_U}(\omega_V)\).  The support of \(J_U\) is locally
set-theoretically principal in \(S_U\) by the geometric
construction.  Lift a local defining function to \(V\) and
invert it.  This exact meromorphic localization realizes
\(j_*\) on the underlying module and gives the description
of \(A_{0,U}\) in the statement.

The two descriptions must agree as inclusions, with the fixed
residue normalization.  On the smooth locus of
\(S_U\setminus J_U\), filtered duality and smooth graph
pushforward give the ordinary dualizing sheaf and its residue
inclusion.  This is precisely Ext-to-support with the
normalization specified by ambient adjunction.  Naturality in
smooth etale coordinates gives the same scalar on each chart.
The agreement extends throughout \(S_U\setminus J_U\), because
the first support-cohomology module has no sections supported
on a smaller-dimensional subset.  One can see this from the
Cousin resolution of the smooth dualizing sheaf: its first term
supported on \(S_U\) is a sum of injective modules at the
codimension-\(c\) generic points.  Localization extends the
agreement across \(J_U\).  We have therefore identified the
actual lowest-piece inclusion, including its residue convention.

We have now identified \(F_0A_{0,U}\) and its inclusion before
direct image.  To identify \(F_0M_U\), compute \(g_+\) by a
graph embedding followed by a smooth
projection.  At filtration zero, only the top term of the relative
Spencer complex survives: it is the direct image of
\(\mathcal B|_{S_U}\) on the graph.  Projective strictness yields
\[
 F_{<0}M_U=0,\qquad
 F_0M_U=R^1g_*(\mathcal B|_{S_U}),
\]
and identifies this sheaf as a coherent subsheaf of the unfiltered
degree-one direct image.  All constructions in this identification
are canonical under etale changes of charts.
\end{proof}

The obstruction sheaf for the transition indexed by \(j,k\) is thus
\(F_0M\otimes C^{-(a+j+k)}\), by
Equation~\eqref{hg:geometric-lines} and projection formula.
We do not need all its global sections to vanish. The lifting proof
will show that the coordinate obstruction lies in the kernel of a
symbol map, and will apply the following vanishing to that kernel.

\subsection{Negative-twist vanishing}

The support of the system \(M\) just constructed is finite over
\(\mathcal P\). We prove the required vanishing for any system
with this support property.

\begin{lemma}\label{hg:negative-vanishing}
Let \(M\) be a compatible system of mixed Hodge modules in perverse
degree zero on the scheme etale charts of \(\mathcal Y\).  Suppose
its support is finite over \(\mathcal P\).  Then, for every integer
\(k\), every positive integer \(b\), and every \(j<0\),
\[
 \mathbb H^j\!\left(
   \mathcal Y,\Gr^F_k\DR_{\mathcal Y}(M)\otimes p^*C^{-b}
 \right)=0.
\]
Here the graded de Rham complex is the descended coherent complex.
\end{lemma}

\begin{proof}
The line \(C\) lives on the root gerbe, so we first push the
complex to \(\mathcal P\).  We will then pull back to projective
space, where \(C\) becomes \(\OO(1)\) and ordinary
Kodaira--Saito vanishing applies.  Since \(p\) is finite on the
support of \(M\), its direct image is concentrated in perverse
degree zero; call the resulting system on \(\mathcal P\) \(M'\).
On ordinary scheme
charts, projective strictness gives
\begin{equation}\label{hg:graded-direct-image}
 Rp_*\Gr^F_k\DR_{\mathcal Y}(M)
 \simeq \Gr^F_k\DR_{\mathcal P}(M').
\end{equation}
For the desired global hypercohomology vanishing, this identity must
hold for the descended complexes.  We verify that compatibility
before passing to a finite cover.

For right differential modules, use the transfer bimodule
\[
 \mathcal D_{\mathcal Y\to\mathcal P}
 =
 \OO_{\mathcal Y}\otimes_{p^{-1}\OO_{\mathcal P}}
                         p^{-1}\mathcal D_{\mathcal P},
\]
with its order filtration.  Derived tensor over
\(\mathcal D_{\mathcal Y}\), followed by derived direct image,
forms the filtered direct image; equivalently, these operations
can be carried out on Rees modules.  De Rham on the base is a
further derived tensor with \(\OO_{\mathcal P}\), filtered from
degree zero.  The filtered Spencer resolution, associativity of
derived tensor, and projection formula identify this with direct
image of the de Rham complex upstairs.  Indeed, tensoring the
transfer module over the base differential operators with the
base structure sheaf gives \(\OO_{\mathcal Y}\).

These sheaf constructions are canonical on the etale site, and the
Spencer resolutions are locally free over differential operators.
They commute with chart changes already before taking associated
gradeds.  The ordinary projective direct-image theorem on the base
charts supplies strictness and concentration in perverse degree
zero.  Passing to the etale site does not change coherent
cohomology on those charts.  This proves
Equation~\eqref{hg:graded-direct-image} globally, as an identity
computing coherent hypercohomology.  The same verification will
apply to the finite representable map below.

To pass from the gerbe to a scheme, use the representable finite
surjective morphism
\[
 v:\PP^s_z\longrightarrow\mathcal P
\]
defined by the power map
\([z_0:\cdots:z_s]\mapsto[z_0^\ell:\cdots:z_s^\ell]\)
and the root \(\OO_{\PP^s_z}(1)\); thus
\(v^*C=\OO_{\PP^s_z}(1)\).  Trivial roots on the standard affine
opens of \(\PP^s\) give scheme etale charts
\(\check U_i\to\mathcal P\).  The restriction of \(v\) to
\(z_i\ne0\) factors through \(\check U_i\).  On each such
chart take the derived Hodge-module inverse image of \(M'\)
and then its degree-zero perverse cohomology.  Thus the object
being constructed chartwise is
\[
 H={}^pH^0(v^*M').
\]
The chosen root and functoriality identify these objects on
overlaps, producing a system on the Zariski opens of \(\PP^s_z\).

We check that \(H\) is a global graded-polarizable mixed Hodge
module, so that ordinary vanishing applies to it.  The perverse,
filtered, and weight data glue, as do the strict-support
decompositions of the pure weight gradeds, by uniqueness.  On a
smooth connected dense stratum of an irreducible support, a
polarization on a nonempty Zariski open extends to a flat pairing:
the fundamental group of that open surjects onto the fundamental
group of the stratum.  Hodge compatibility extends by continuity,
and the extended flat pairing remains nondegenerate and positive.
The local Hodge modules already give boundary quasi-unipotence.
Saito's extension theorem for polarizable variations
\cite[Section~3.b]{Saito1990} and uniqueness of strict-support
extension identify the global pure object with the glued object.
Projectivity of \(\PP^s_z\) leaves no additional boundary at
infinity.  The weight extensions give the required mixed object.

Vanishing for \(H\) will imply vanishing for \(M'\) once we
realize \(M'\) as a retract of \(v_+H\).  On a chart one
must use the \emph{whole} base-changed finite cover.  Over
\(\check U_i\), this is a disjoint union of copies of the
corresponding affine coordinate chart of \(\PP^s_z\).  The unit
on unshifted constants and its dual trace, formed using smooth
duality in equal dimensions, compose to multiplication by the
degree.  This can first be checked on the finite etale locus,
where the trace sums over the sheets.  On each connected smooth
base chart, the shifted constant Hodge module is the rank-one
intersection complex, with endomorphism ring \(\Q\), and its
endomorphisms are determined on a dense open.  The composition
therefore equals the degree globally, with no contribution
supported on the branch locus.  This argument permits
ramification.

Tensor the unit and trace with \(M'\), apply proper projection
formula, and take perverse cohomology in degree zero.  Finite
direct image is perverse exact, so it commutes with this
degree-zero cohomology.  We obtain maps
\[
 M'\longrightarrow v_+H\longrightarrow M'
\]
whose composite is multiplication by the degree of the whole
base-changed cover.  Dividing the second map by that degree
gives the retraction.  Each map is canonical under etale base change; the
retraction consequently descends on filtered differential modules
as well.

Applying the finite version of
Equation~\eqref{hg:graded-direct-image} and projection formula
now exhibits the hypercohomology in the statement as a direct
summand of
\[
 \mathbb H^j\!\left(
   \PP^s_z,\Gr^F_k\DR(H)\otimes\OO_{\PP^s_z}(-b)
 \right).
\]
Ordinary negative-ample Kodaira--Saito vanishing makes this group
zero for \(j<0\).  For a mixed object, apply the vanishing to
the pure weight gradeds and then to the weight-filtration exact
sequences, whose Hodge filtrations are strict.
\end{proof}

\subsection{Lifting through the infinitesimal neighborhoods}

We now combine the lowest-piece identification with negative-twist
vanishing. The remaining step is to identify the relation satisfied
by the coordinate and conormal obstructions.

\begin{proof}[Proof of Proposition~\ref{prop:formal-lifting}]
We induct on \(k\), simultaneously for all powers \(I^j\).
The connecting homomorphisms will factor through the first
graded layers and together form a derivation.  We must show
that this derivation vanishes both on functions from \(U\)
and on a local generator of \(I/I^2\).  The graph calculation
relates these two errors in the differential module \(M_U\).
Its order-one symbol, together with negative-twist vanishing,
first kills the function error.  The original relation then
kills the conormal error.

Define
\[
 E_{j,k}=g_*(I^j/I^{j+k}),\qquad j\geq0,\quad k\geq1.
\]
The first graded layer is
\[
 E_{j,1}=\OO_{\mathcal T_U}\otimes C^{-j}.
\]
For adjacent lengths, the exact sequence has kernel
\(g_*\OO_{S_U}\otimes C^{-j-k}\), and its connecting homomorphism
takes values in
\[
 R^1g_*\OO_{S_U}\otimes C^{-j-k}.
\]

\smallskip\noindent\emph{The obstruction derivation.}
Assume inductively that all shorter transitions are surjective.
A section of
\(E_{j,k}\) with zero length-one term comes from \(E_{j+1,k-1}\)
when \(k>1\).  Induction lifts it locally from \(E_{j+1,k}\),
so its connecting class vanishes.  Induction also makes
\(E_{j,k}\to E_{j,1}\) surjective.  The obstruction thus
factors through the length-one layer; for \(k=1\) this
factorization is immediate.  More precisely, these factorizations
give additive maps
\[
 D_{k,j}:\OO_{\mathcal T_U}\otimes C^{-j}
       \longrightarrow R^1g_*\OO_{S_U}\otimes C^{-j-k}.
\]
Together they define \(D_k\) on the graded algebra
\(\bigoplus_{j\geq0}\OO_{\mathcal T_U}\otimes C^{-j}\), with
graded shift \(k\); no \(\OO_U\)-linearity is asserted.
Its values are computed by choosing local
lifts and taking Cech differences.  For a product, this
difference is the sum of the two first-order differences:
the product of two errors vanishes in the layer under
consideration.  Thus \(D_k\) is a \(\C\)-derivation, with
values in the graded module whose degree-\(j\) term is
\(R^1g_*\OO_{S_U}\otimes C^{-j}\).

Restrict its degree-zero part along
\(\OO_U\to\OO_{\mathcal T_U}\).  This derivation factors
through \(\Omega_U^1\), defining a section of
\(T_U\otimes R^1g_*\OO_{S_U}\otimes C^{-k}\).  As the chart
varies, these sections descend to
\begin{equation}\label{hg:obstruction-section}
 e\in H^0\!\left(
 \mathcal Y,T_{\mathcal Y}\otimes R^1g_*\OO_S\otimes C^{-k}
 \right)
 =
 H^0\!\left(
 \mathcal Y,T_{\mathcal Y}\otimes F_0M\otimes C^{-(a+k)}
 \right).
\end{equation}
The equality follows from Equation~\eqref{hg:geometric-lines},
Lemma~\ref{hg:lowest-piece}, and projection formula.  For the
descent assertion, function lifts pull back along the unique
etale thickenings.  Ordinary flat base change on the reductions
pulls back their Cech classes in the coherent obstruction
sheaves.  Differentials also pull back and generate under an
etale map, giving the required compatibility.

\smallskip\noindent\emph{Coordinate errors and adjunction.}
We now compute this derivation in local coordinates.  After shrinking \(U\), choose etale coordinates
\(t_1,\ldots,t_d\), with \(d=\dim U\), and a frame of
\(C^{-1}\).  Write \(y\) for its pullback, a conormal frame.
Induction lifts the coordinates along \(g\) modulo \(I^k\)
and lifts \(y\) to \(I/I^{k+1}\).  Lift them one step further
on an open cover of \(S_U\), obtaining functions
\(h_{i,\lambda}\) modulo \(I^{k+1}\) and generators \(y_i\)
modulo \(I^{k+2}\).  The conormal generator must be lifted one
order farther than the coordinates: its next error lies in
\(I^{k+1}/I^{k+2}\), whereas a coordinate error lies in
\(I^k/I^{k+1}\).  By formal etaleness of the coordinate map,
each \(h_i\) defines a local map from \((k+1)S_U\) to \(U\).
Write the overlap errors as
\begin{equation}\label{hg:coordinate-errors}
 h_{j,\lambda}-h_{i,\lambda}
       =e_{ij,\lambda}y_i^k,
 \qquad
 y_j-y_i=\eta_{ij}y_i^{k+1}.
\end{equation}
Thus \(D_k(t_\lambda)\) is represented by
\((e_{ij,\lambda}y^k)\), and \(D_k(y)\) by
\((\eta_{ij}y^{k+1})\).  The powers of \(y\) record the
target line bundles; \(e_{ij,\lambda}\) and \(\eta_{ij}\)
are their coefficients on the reduction.  Denote the induced
frame of
\(\mathcal B|_{S_U}\) by \(\sigma=y^{-a}\).

Off \(J_U\), the fixed adjunction isomorphism gives
\[
 \omega_{(k+1)S_U}
 \simeq\OO_{\mathcal X}((a+k)S)|_{(k+1)S_U}.
\]
Via this isomorphism, the local equation \(y_i\) defines a
dualizing generator, denoted \(b_i=y_i^{-(a+k)}\).  This
notation denotes a frame of the line bundle
\(\OO_{\mathcal X}((a+k)S)\) restricted to the thickening;
it does not invert the nilpotent function \(y_i\) in its
structure sheaf.  Realize the dualizing sheaves in support cohomology.  Under
the trace inclusion
\(\omega_{S_U}\hookrightarrow\omega_{(k+1)S_U}\), we have
\begin{equation}\label{hg:dualizing-errors}
 b_j-b_i=-(a+k)\eta_{ij}\sigma,\qquad
 y_i^kb_i=\sigma,\qquad y_i^{k+1}b_i=0.
\end{equation}
The first identity is the binomial expansion of
\((1+\eta_{ij}y_i^k)^{-(a+k)}\): all quadratic terms vanish
because \(2k\geq k+1\).  In the last two identities,
\(\sigma\) is viewed in the dualizing sheaf of the thickening
through the trace inclusion.  They express Cartier trace and
the annihilator of the thickening.

\smallskip\noindent\emph{The graph calculation.}
We next put both errors in the same differential module.  Embed
\((k+1)S_U\) as a closed subscheme of a smooth open
\(Z\subset\overline Z=\PP^m\).  The graph of the reduced map
\(g\) is closed in \(\overline Z\times U\), since it is proper
over \(U\).  Let \(\iota:S_U\hookrightarrow\overline Z\times U\)
be this reduced graph embedding, and let \(\mathcal N_\Gamma\)
be the underlying unfiltered differential module of
\(\iota_+A_{0,U}\).  This module is distinct from the direct-image
module \(M_U\) on \(U\).  For each local graph lift \(h_i\),
Ext-to-support cohomology sends the dualizing section \(b_i\)
to a section \(\delta_i(b_i)\) of \(\mathcal N_\Gamma\).
All the lifted graphs have the same reduced support, namely
the graph of \(g\), so these sections lie in one support-cohomology
module rather than in different modules for the different lifts.
These sections need not lie in its lowest filtered piece.
We will prove the identity
\begin{equation}\label{hg:graph-identity}
 \delta_j(b_j)-\delta_i(b_i)
 =
 -(a+k)\eta_{ij}\sigma
 +\sum_{\lambda=1}^d
          (e_{ij,\lambda}\sigma)\cdot\partial_{t_\lambda}.
\end{equation}
On the right, each reduced dualizing section \(\sigma\) is
included in \(\mathcal N_\Gamma\) by the lowest-piece map of
Lemma~\ref{hg:lowest-piece}.  The right differential-operator
action is taken in this graph module.

The calculation also applies when \(S_U\) is singular.  Put
\(c_0=\dim Z-N\).  Ext-to-support identifies the dualizing
sheaf of the thickening with its annihilator submodule in
\(\mathcal H^{c_0}_{S_U}(\omega_Z)\).  To see this, use the
support-cohomology spectral sequence.  Support cohomology
vanishes below \(c_0\), so in total degree \(c_0\) the Ext
group consists of homomorphisms from the structure sheaf of
the thickening to the first support-cohomology module.  In
particular, no local complete-intersection hypothesis on
\(S_U\subset Z\) is needed.

Lift the \(h_{i,\lambda}\) locally to functions on \(Z\).
For a dualizing section \(b\) of the thickening, its graph
inclusion in \(Z\times U\) is the generalized fraction
\[
 \frac{b\,dt_1\wedge\cdots\wedge dt_d}
      {\prod_{\lambda=1}^d(t_\lambda-h_{i,\lambda})}.
\]
It represents successive support cohomology in the additional
coordinate equations, equivalently the Koszul residue for the
graph, and can be computed etale-locally near the graph branch.
After taking support cohomology from \(Z\), the translated
new coordinates form a regular sequence; their support
cohomology is consequently concentrated in the degree equal
to their number.  Either set of translated coordinates gives
the same localization of this support-torsion module, because
the two translations differ nilpotently on every section.

The change from \(h_i\) to \(h_j\) thus has a finite Taylor
expansion on \(b_j\).  Again using \(2k\geq k+1\), products
of two differences from Equation~\eqref{hg:coordinate-errors}
annihilate \(b_j\).  The linear terms are determined by
\[
 (h_{j,\lambda}-h_{i,\lambda})b_j
       =e_{ij,\lambda}\sigma.
\]
The right action on a top differential form satisfies
\[
 \left(\frac{dt_\lambda}{t_\lambda-h}\right)
       \cdot\partial_{t_\lambda}
   =\frac{dt_\lambda}{(t_\lambda-h)^2}.
\]
The doubled pole therefore contributes the positive sign in
Equation~\eqref{hg:graph-identity}.  Combining these terms
with Equation~\eqref{hg:dualizing-errors} proves the identity.
So far the calculation is off \(J_U\).  Both its sections
and its identities extend meromorphically across \(J_U\) in
the \emph{unfiltered} localized module, where arbitrary finite
pole orders are allowed.

\smallskip\noindent\emph{Vanishing of the obstruction.}
Push the graph identity along the projection
\(\pi:\overline Z\times U\to U\), using the relative Spencer
complex.  Closed direct image is perverse exact, and
\(\pi_+\iota_+A_{0,U}=g_+A_{0,U}\).  Thus degree-one holonomic
differential-module cohomology of this Spencer direct image is
the underlying module of
\(M_U={}^{p}\!H^1g_*A_{0,U}\).
The cochain \((\delta_i(b_i))\) belongs to the top, degree-zero
term, from which there is no outgoing relative Spencer
differential.  Its Cech difference is a boundary in total
degree one.  A sufficiently refined ambient cover around
the closed graph computes this boundary; all the sheaves
in question are supported there.  By
Lemma~\ref{hg:lowest-piece}, the reduced cocycles in
Equation~\eqref{hg:graph-identity} represent their
\(R^1g_*\mathcal B\) classes in \(F_0M\).  Right
differentiation in the \(U\)-coordinates acts on the
pushforward complex.  Consequently the Cech classes satisfy
\[
 -(a+k)[\eta_{ij}\sigma]
 +\sum_{\lambda=1}^d
       [e_{ij,\lambda}\sigma]\cdot\partial_{t_\lambda}=0
 \qquad\text{in }M_U.
\]
Here each bracketed class belongs to
\(R^1g_*\mathcal B=F_0M_U\); after differentiation, its
image belongs to \(F_1M_U\).  Thus every term of the relation
lies in \(F_1M\), and the \(\eta\)-term lies in \(F_0M\).
The equality itself was proved in the unfiltered differential
module.  The inclusion \(F_1M\subset M\)
is an inclusion of actual sheaves, so vanishing in \(M\)
implies vanishing in \(F_1M\).  Taking the order-one symbol
shows that the global section \(e\) from
Equation~\eqref{hg:obstruction-section} is killed by
\begin{equation}\label{hg:symbol-map}
 T_{\mathcal Y}\otimes F_0M
           \longrightarrow \Gr^F_1M,
\end{equation}
after twisting by \(C^{-(a+k)}\).
Only the relation is used in this filtered step; no bound
on the Hodge filtration of the auxiliary cochain \((\delta_i(b_i))\)
has been imposed.

Because \(F_{<0}M=0\), the entire complex at filtration
index one is
\[
 \Gr^F_1\DR(M)=
 \bigl[T_{\mathcal Y}\otimes F_0M
       \longrightarrow\Gr^F_1M\bigr],
\]
in degrees \(-1\) and \(0\), with differential the symbol
map in Equation~\eqref{hg:symbol-map}.  After the indicated
twist, its degree-\(-1\) hypercohomology is therefore exactly
the space of global sections of the kernel of that map.
The twist is negative,
since \(a+k>0\).  Lemma~\ref{hg:negative-vanishing}
annihilates this space, and hence \(e=0\).

It remains to kill the conormal part of the derivation.
With coordinates and a line trivialization fixed locally,
\(e=0\) says that every class
\([e_{ij,\lambda}\sigma]\) vanishes in \(F_0M\), and thus
in \(M\).  Its actual differential-operator image vanishes
as well, before passing to symbols.  Returning to the
unfiltered relation gives
\[
 (a+k)[\eta_{ij}\sigma]=0.
\]
Injectivity of the lowest-piece inclusion and \(a+k>0\)
then give \(D_k(y)=0\).  The degree-zero derivation also
vanishes: its restriction to \(\OO_U\) is the coordinate
error \(e\), and \(\OO_U\to\OO_{\mathcal T_U}\) is
surjective because \(\mathcal T_U\) is a closed subscheme
of \(U\).  The coefficient
algebra and the conormal frame \(y\) generate the full
graded length-one algebra locally.  Thus \(D_k=0\) in
every degree, completing the induction and proving
surjectivity of all the sheaf transitions.

The kernels are the stated coherent sheaves
\(\OO_{\mathcal T_U}\otimes C^{-j-k}\).  If \(U\) is
affine, their first cohomology vanishes, so the exact
sequences of sheaves of abelian groups give surjectivity
on global sections as well.  Successive choices of lifts
then give compatible formal lifts of any chosen functions
on \(\mathcal T_U\) and, after trivializing \(C\), of the
transverse conormal frame.
\end{proof}
